\begin{table}[htbp]
\centering
\caption{Within-model frame-manipulation effects}
\label{tab:h3_frame_manipulation}
\begin{tabular}{lrrrrr}
\toprule
Frame contrast & Coefficient & SE & \(t\) & 95\% CI & \(p\) \\
\midrule
Autonomy vs. neutral & 2.6204 & 0.3440 & 7.618 & [1.9458, 3.2949] & 0.0000 \\
Collective obligation vs. neutral & -10.5426 & 0.3440 & -30.649 & [-11.2172, -9.8680] & 0.0000 \\
Equity/access vs. neutral & 6.2870 & 0.3440 & 18.277 & [5.6125, 6.9616] & 0.0000 \\
\midrule
\(N\) & \multicolumn{5}{r}{2160} \\
\(R^2\) & \multicolumn{5}{r}{0.6653} \\
Adjusted \(R^2\) & \multicolumn{5}{r}{0.6603} \\
\bottomrule
\end{tabular}
\begin{flushleft}
\footnotesize Notes: The dependent variable is the profile--repetition-level burden effect, \(\Delta = W^{low} - W^{high}\). The neutral frame is the omitted reference category. The pooled model includes four base models: GPT-4.1, Llama 3.1 70B Instruct, Qwen3.6-72B Instruct, and Mistral Large, with model and profile fixed effects. Positive coefficients indicate larger predicted reductions in willingness under high administrative burden.
\end{flushleft}
\end{table}
